\chapter{Unifying Decoding to a Single Greedy Path (Phase 1.6)}
\label{chap:p1-6}

\section{Goal}

The pilot evaluated its metrics across three different decoding regimes without treating that as a
variable. Baseline inference, descriptive-accuracy masking, and robustness all used beam search
(\texttt{num\_beams=3}); attribution used greedy decoding (\texttt{num\_beams=1}); stability used
sampling (\texttt{do\_sample=True}). Because decoding is not a neutral implementation detail --- it
changes which token sequence, and therefore which grounded box, the model emits --- metrics computed on
different decoding paths are not strictly comparable, and, worse, an explanation computed on one path
can be explaining a different answer than the baseline reported on another.

The plan (item 3.6) prescribes adopting greedy decoding as the single canonical deterministic path, for
two stated reasons: it is the path attribution can be extracted from without beam-reordering
bookkeeping, and the pilot's own efficiency profiling suggested beam and greedy cost about the same for
short grounding outputs, so nothing is lost by switching. Stability's sampling stays (that stochasticity
is the metric's point) but is anchored against the greedy path.

\section{What Was Done}

The unification is centralized rather than scattered. \texttt{config.DECODING} is a single switch
(\texttt{"beam"} $\rightarrow$ 3 beams, the frozen default; \texttt{"greedy"} $\rightarrow$ 1), and
\texttt{model.run\_task} now resolves its beam count from that switch when no explicit value is passed
(\texttt{\_effective\_num\_beams}). Because every deterministic grounding call --- baseline, masking,
robustness, sparsity's region lookups --- flows through \texttt{run\_task} without specifying beams,
flipping one config value moves all of them onto greedy together. The two paths that legitimately
differ pass their beam count explicitly and are therefore untouched: stability
(\texttt{do\_sample=True, num\_beams=1}) and attribution (its own greedy \texttt{generate}). Sparsity
was already greedy, since it reads cross-attention from the attribution path.

The default is \texttt{"beam"}, so the frozen pilot reproduces byte-for-byte; scripts 03, 05, and 09
expose a \texttt{-{}-decoding} flag, and any non-default run writes to a \texttt{greedy}-suffixed file.
Four unit tests pin the resolver: explicit values always win, and \texttt{None} resolves from the
current policy.

\section{Results}

Re-running the full 163-sample baseline under greedy and comparing to the frozen beam baseline:

\begin{center}
\begin{tabular}{p{5.8cm}p{4.0cm}}
\toprule
\textbf{measure} & \textbf{beam vs.\ greedy} \\
\midrule
answer agreement (overall)      & 93.3\% (11 of 163 differ) \\
\ \ rule\_1 / rule\_2 / rule\_3 / rule\_4 & 92.1\% / 92.3\% / 93.9\% / 96.0\% \\
mean inference time             & 271 ms (beam) vs.\ 261 ms (greedy), ratio 0.96 \\
median inference time           & 261 ms vs.\ 250 ms \\
\bottomrule
\end{tabular}
\captionof{table}{Greedy vs.\ beam baseline over 163 samples. Greedy costs essentially the same
(0.96$\times$), confirming the efficiency rationale, and agrees with beam on 93.3\% of answers.}
\label{tab:p16-agree}
\end{center}

Both plan claims hold. Greedy is not more expensive --- it is marginally cheaper (0.96$\times$),
consistent with the pilot's observation that beam search buys little for outputs this short. And the
answer agreement is high, 93.3\%, so unifying onto greedy does not upend the pilot's conclusions.

The more important number is the 6.7\% that \emph{disagree}, because it quantifies a fragmentation the
pilot could not see. Under the old regime, the baseline reported a beam-search answer while the
attribution (sparsity, heatmaps) was computed under greedy. For roughly one sample in fifteen, those
two paths produce different grounded boxes and therefore potentially different answers --- meaning the
pilot's own explanation for that sample was, strictly, explaining a decoding path the baseline never
reported. Unifying the deterministic path removes that mismatch by construction: after the switch, the
answer a metric reports and the answer its attribution explains are produced by the same decoder. The
93.3\% agreement is thus the measure of how much the pilot's cross-metric decoding inconsistency
mattered --- modest in aggregate, real at the per-sample level, and now eliminated.

The frozen \texttt{baseline\_predictions.csv} is untouched (empty \texttt{git diff}); the greedy run
wrote \texttt{baseline\_predictions\_greedy.csv}. The full suite is green at 109 tests.

\section{Standard vs. Non-Standard Choices}

\textbf{Standard.} Fixing one decoding configuration across an evaluation, rather than letting each
metric inherit whatever default its code happened to use, is basic experimental control.

\textbf{Non-standard, and the contribution.} Two things. The unification is a \emph{single} switch read
at call time, so there is exactly one place decoding is decided and no metric can silently drift onto a
different path --- and the switch defaults to the frozen policy, so the pilot's beam-search numbers
remain the reproducible baseline. And the choice of greedy is justified by evidence rather than
convention: it was adopted specifically because it is the path attribution can be read from, and only
after confirming --- here, at 0.96$\times$ --- that it costs no more than the beam search it replaces.
The plan flags one caveat this validates conditionally: if a later decision adopts dense captioning
(long outputs), the beam-vs-greedy cost equivalence must be re-measured, because the decode loop stops
being nearly free.

\section{Implications}

Attribution (Phase 4) can now be read from the same decoder that produces the reported answer, so a
saliency map or cross-attention heatmap explains the answer the baseline actually gave --- a
precondition for the two-stream attribution-agreement work. Every deterministic metric shares one
execution path, so cross-metric comparisons are no longer confounded by decoding. And the switch is one
config line for the scale run, with the frozen beam numbers preserved beside the greedy ones for
reference.
